\documentclass[12pt, a4paper]{article}
\usepackage[utf8]{inputenc}
\usepackage{geometry}
\geometry{a4paper, margin=1in}
\usepackage{hyperref}
\usepackage{titlesec}
\usepackage{setspace}
\usepackage{graphicx}
\usepackage{caption}
\usepackage{tabularx}
\usepackage{booktabs}
\usepackage{longtable}
\usepackage{pdflscape}

\title{\textbf{Edge AI Surveillance Systems Using TinyML on ESP32 Microcontrollers with CNN Quantization for Privacy-Preserving Local Processing: A Literature Review}}
\author{}
\date{}

\begin{document}

\maketitle

\begin{abstract}
Evidence on edge AI surveillance systems using TinyML on ESP32 microcontrollers with CNN quantization for privacy-preserving local processing remains limited, with no studies directly deploying quantized CNNs on ESP32 hardware for surveillance tasks; instead, findings demonstrate feasibility on comparable low-power microcontrollers like MAX78000 (up to 180 fps, 196 µJ per inference) and general edge devices achieving 16 fps with 89\% precision \cite{moosmann2023, zhao2020}. Privacy-preserving local processing is consistently supported through edge-based anonymization of faces and minors (92.1\% classification accuracy), semantic scene graph generation, and chaotic frame scrambling, reducing transmission risks compared to cloud models \cite{yuan2020, huang2023, fitwi2021}. Lightweight CNN optimizations, including depthwise separable convolutions, adaptive 4-bit quantization, and quantized TinyML networks ($<$0.5 MB memory), enable real-time object detection (e.g., handgun/knife, human tracking) on resource-constrained devices, with strong emphasis on latency and bandwidth reduction via edge-cloud hybrids \cite{zhao2020, chin2022, moosmann2023}. This review synthesizes advancements in TinyML for surveillance amid rising smart city demands, addressing gaps in microcontroller-specific deployments. Key implications include scalable privacy protection without central data aggregation, though hardware specificity to ESP32 is absent, highlighting needs for targeted benchmarks. Future work should prioritize ESP32 validations to realize full local processing potential.

\textbf{Keywords:} edge, computing, surveillance, privacy, preserving, latency, bandwidth, reduction
\end{abstract}

\section{Introduction}
The proliferation of video surveillance in smart cities, campuses, and industrial sites has transformed monitoring into a cornerstone of public safety, situational awareness, and operational efficiency. Traditional centralized systems, reliant on cloud processing of raw video streams, face critical bottlenecks: high latency from data transmission, massive bandwidth demands, and privacy risks from exposing personally identifiable information such as faces or minors in transit \cite{myneni2022, yuan2020}. Edge AI, particularly TinyML on ultra-low-power microcontrollers, emerges as a paradigm shift by enabling local inference with convolutional neural networks (CNNs), quantized to fit severe memory constraints (e.g., $<$0.5 MB), thereby preserving privacy through on-device processing without raw data upload \cite{moosmann2023}. Quantization techniques, including post-training and quantization-aware methods, compress models to 4-8 bits while targeting real-time performance in milliwatt power domains, critical for battery-operated surveillance nodes.

Despite these promises, integrating TinyML surveillance on specific platforms like ESP32 microcontrollers---known for Wi-Fi integration and $\sim$520 KB SRAM---remains underexplored, with most work on analogous hardware like STM32 or MAX78000. Challenges persist in balancing accuracy, inference speed (e.g., FPS), energy efficiency, and privacy for tasks like object detection amid heterogeneous edge environments. This review synthesizes evidence on edge AI surveillance systems using TinyML on microcontrollers with CNN quantization, focusing on privacy-preserving local processing, to clarify feasibility, optimizations, and gaps relative to the research question: ``Edge AI Surveillance Systems using TinyML on ESP32 microcontrollers with CNN quantization for privacy-preserving local processing.''

\section{Methods}
\subsection{Search Strategy}
We performed a comprehensive search across over 220 million academic papers from Semantic Scholar and OpenAlex databases. The search strategy employed hybrid semantic and keyword-based retrieval to maximize coverage.

Search queries included:
\begin{itemize}
    \item ``edge-computing surveillance IoT edge-AI privacy-preserving low-latency bandwidth-reduction CCTV''
    \item ``TinyML microcontroller computer-vision CNN edge-device image-classification detection''
    \item ``CNN-quantization post-training-quantization quantization-aware-training model-compression TinyML edge-AI''
    \item ``ESP32 TinyML surveillance object-detection edge-AI microcontroller vision''
    \item ``cloud-to-edge surveillance transition intelligent-video-analytics edge-computing AI-CCTV''
    \item ``IoT-edge surveillance real-time privacy low-latency bandwidth-efficiency video-analytics''
    \item ``microcontroller-CNN challenges TinyML memory-constrained inference computer-vision''
\end{itemize}

\subsection{Study Selection}
Initial database searching identified 280 records. After duplicate removal and relevance-based filtering, 67 records were screened against eligibility criteria. Of these, 47 papers were excluded, resulting in 20 papers included in the final synthesis.

\textbf{PRISMA Flow Diagram Summary:}
\begin{itemize}
    \item Records identified through database searching: n = 280
    \item Records after duplicates removed and filtered: n = 67
    \item Records screened: n = 67
    \item Records excluded: n = 47
    \item Full-text articles assessed for eligibility: n = 20
    \item Studies included in qualitative synthesis: n = 20
\end{itemize}

\textbf{Eligibility criteria included:}
\begin{itemize}
    \item \textbf{Edge AI Focus:} Does the paper address edge computing or TinyML for AI inference on resource-constrained devices (not cloud-only)?
    \item \textbf{Surveillance Application:} Is the work applied to surveillance, video analytics, object detection, or computer vision tasks relevant to monitoring?
    \item \textbf{Microcontroller Deployment:} Does the paper deploy models on microcontrollers like ESP32, STM32, or similar (not just Raspberry Pi/Jetson Nano)?
    \item \textbf{Model Optimization:} Does it discuss CNN quantization, pruning, compression, or TinyML optimizations for vision models?
    \item \textbf{Privacy or Edge Benefits:} Does it mention benefits like privacy preservation, low latency, or bandwidth reduction from edge processing?
    \item \textbf{Real-World Evaluation:} Does it include experimental results on actual hardware with metrics like FPS, power, accuracy?
    \item \textbf{Recent Publication:} Published 2019 or later?
\end{itemize}

\textit{Note on included studies:} Multiple papers were retained despite not meeting all criteria strictly, particularly the "Microcontroller Deployment" requiring ESP32 or similar and "Real-World Evaluation" demanding specific metrics like FPS/power on hardware. For example, \cite{myneni2022, zhao2020, liu2023}, and others emphasize edge architectures without naming microcontrollers like ESP32/STM32 or providing exact FPS/power metrics (noted as "not reported"), yet were included as they provide the only direct evidence on privacy-preserving surveillance frameworks and lightweight CNNs directly relevant to TinyML contexts. Similarly, general TinyML papers like \cite{aatab2023} and \cite{roveri2022} lack surveillance tasks but offer foundational on-device training evidence absent elsewhere. These exceptions maintain transparency and comprehensiveness for synthesizing the narrow research question.

\subsection{Data Extraction and Synthesis}
Data extraction focused on the following variables:
\begin{itemize}
    \item \textbf{Key Contributions:} Summarize main innovations in edge AI surveillance, TinyML deployment, or optimization techniques.
    \item \textbf{Hardware Platform:} Specific microcontroller or edge device used (e.g., ESP32) with key constraints like RAM, flash, power.
    \item \textbf{Optimization Methods:} Techniques applied: quantization (PTQ/QAT), pruning, knowledge distillation, architecture search.
    \item \textbf{Performance Results:} Metrics: inference time (ms), accuracy (\%), memory footprint (KB), power (mW), FPS.
    \item \textbf{Surveillance Tasks:} Applications: object detection, classification, anomaly detection, privacy features.
    \item \textbf{Limitations:} Reported challenges: accuracy loss, real-time limits, deployment issues.
    \item \textbf{Publication Details:} Year, venue, authors (first few).
\end{itemize}
Thematic analysis was employed to identify patterns and synthesize findings across studies. Evidence strength was assessed based on consistency of findings and number of supporting studies.

\section{Results}
\subsection{Characteristics of Included Studies}

\begin{landscape}
\footnotesize
\begin{longtable}{ p{3.5cm} p{0.8cm} p{3.5cm} p{3cm} p{3.5cm} p{3.5cm} p{3.5cm} }
\caption{Characteristics of Included Studies} \label{tab:characteristics} \\
\toprule
\textbf{Study ID} & \textbf{Year} & \textbf{Key Focus} & \textbf{Hardware Platform} & \textbf{Optimization Methods} & \textbf{Surveillance Tasks} & \textbf{Performance Results} \\
\midrule
\endfirsthead
\toprule
\textbf{Study ID} & \textbf{Year} & \textbf{Key Focus} & \textbf{Hardware Platform} & \textbf{Optimization Methods} & \textbf{Surveillance Tasks} & \textbf{Performance Results} \\
\midrule
\endhead
\midrule
\multicolumn{7}{r}{\textit{Continued on next page}} \\
\endfoot
\bottomrule
\endlastfoot

(Myneni et al., 2022) \cite{myneni2022} & 2022 & Privacy-preserved smart-city video surveillance (SCVS) & Distributed edge cloud & Parameter server-based distributed ML & Event identification, face anonymization & Not reported \\

(Zhao et al., 2020) \cite{zhao2020} & 2020 & Lightweight deep learning for intelligent edge surveillance & Edge device & Depthwise separable convolutions & Object recognition from surveillance video & 16 fps, 89\% precision \\

(Moosmann et al., 2023) \cite{moosmann2023} & 2023 & TinyissimoYOLO for TinyML object detection & MAX78000, STM32H7A3, STM32L4R9, Apollo4b & 8-bit quantization, QAT & Object detection (up to 3 classes) & 180 fps, 196 µJ/inference, 422k params \\

(Yuan et al., 2020) \cite{yuan2020} & 2020 & Minor privacy protection via edge video processing & Edge surveillance system & Cascaded deep learning & Minor detection, face anonymization & 92.1\% classification accuracy \\

(Aatab \& Freitag, 2023) \cite{aatab2023} & 2023 & On-device training library for TinyML & Microcontroller boards & None specified & Not specified & Not reported \\

(Chin et al., 2022) \cite{chin2022} & 2022 & Adaptive quantization for CNN inference & Resource-constrained edge & Adaptive 4-bit INT4 quantization & Not specified & Better than state-of-the-art 4-bit models \\

(Liu \& HU, 2023) \cite{liu2023} & 2023 & Edge handgun/knife detection & Edge devices in IoT & Single-stage CNN architecture & Handgun/knife detection & Not reported \\

(Nikouei et al., 2018) \cite{nikouei2018sec} & 2018 & Smart surveillance for human detection/tracking & Edge-fog layers & Lightweight algorithms & Human detection/tracking & Not reported \\

(Hailesellasie \& Hasan, 2019) \cite{hailesellasie2019} & 2019 & MulNet flexible CNN processor & FPGAs (Xilinx Zynq SoCs) & Fixed-point/power-of-2 quantization & Not specified & 136.9 GOP/s, 88.49 GOP/s/W \\

(Huang et al., 2023) \cite{huang2023} & 2023 & Semantic privacy-preserving video surveillance & Edge devices & Scene graph generation & Video monitoring, human privacy protection & Not reported \\

(Heredia \& Barros-Gavilanes, 2019) \cite{heredia2019} & 2019 & Edge vs. cloud for surveillance image processing & Embedded device with GPU vs. server & Transfer learning, homography & Object detection/counting & 92\% precision (edge), 93\% (server) \\

(Xu et al., 2021) \cite{xu2021} & 2021 & FL-YOLO for coal mine surveillance & NVIDIA Jetson TX1 (edge) & Depthwise separable convolutions & Object detection (pedestrians) & 36.7 fps, 76.7\% AP, 16.1 MB model \\

(Fitwi et al., 2021) \cite{fitwi2021} & 2021 & Lightweight frame scrambling for privacy & Resource-constrained edge & Chaotic methods (DyCIE, RoI-Mask) & Privacy protection in video streams & Not reported \\

(Banbury et al., 2024) \cite{banbury2024} & 2024 & Wake Vision TinyML dataset/benchmark & Low-power TinyML devices & Dataset generation pipeline & Binary classification (person detection) & Up to 6.6\% accuracy improvement \\

(Kim \& Kim, 2020) \cite{kim2020} & 2020 & Mixture quantization for lightweight CNN & Embedded/mobile & Deterministic/stochastic quantization & Image classification/object detection & Up to 0.44\% mAP gain \\

(Gheisari et al., 2019) \cite{gheisari2019} & 2019 & ECA privacy architecture for smart city & IoT edge & Ontology-based privacy model & Privacy-preserving sensing & Not reported \\

(Fagbohungbe et al., 2022) \cite{fagbohungbe2022} & 2022 & Privacy-preserving edge framework for image classification & Edge devices & Autoencoder latent vectors & Image classification & Not reported \\

(Roveri \& Pavan, 2022) \cite{roveri2022} & 2022 & TinyML on-device training toolbox & Tiny hardware MCUs & On-device training for FFNN/CNN & Not specified & Not reported \\

(Lyu \& Luo, 2022) \cite{lyu2022} & 2022 & Edge-YOLO for airport surveillance & Edge-end embedded devices & Model lightweighting, attention mechanism & Object detection & Not reported \\

\end{longtable}
\end{landscape}

\normalsize

Studies span 2018-2024, predominantly conference proceedings (e.g., IEEE AICAS, SEC) and journals (e.g., IEEE Access, Transactions). Focus areas cluster around privacy-preserving edge processing, lightweight CNNs for object detection, and TinyML optimizations, with hardware mostly unspecified or general edge/microcontrollers; specific metrics like FPS/energy are sparsely reported.

\subsection{Thematic Findings}

\subsubsection{Privacy-Preserving Local Processing in Edge Surveillance}
Edge-based anonymization and semantic processing consistently enable privacy protection by avoiding raw video transmission, with approaches including face/minor detection (92.1\% accuracy \cite{yuan2020}), in-memory encryption, and scene graph generation to safeguard human images \cite{myneni2022, huang2023}. Chaotic scrambling methods (DyCIE, RoI-Mask, SCM) provide lightweight end-to-end privacy at the edge, computationally feasible versus cryptographic alternatives \cite{fitwi2021}. These outperform centralized models by preserving data ownership during distributed training \cite{myneni2022}. Confidence: Moderate (consistent across multiple studies with architectural focus).

\subsubsection{CNN Quantization and Lightweight Optimizations for TinyML}
Quantization techniques dominate model compression: 8-bit quantization-aware training for $<$0.5 MB memory (422k parameters \cite{moosmann2023}), adaptive 4-bit INT4 adjusting scaling/shifting for biased activations (surpassing full-precision in cases \cite{chin2022}), and deterministic/stochastic mixes (0.44\% mAP gain classification, 0.91\% detection \cite{kim2020}). Depthwise separable convolutions reduce complexity in INES (16 fps, 89\% precision \cite{zhao2020}) and FL-YOLO (36.7 fps, 76.7\% AP, 16.1 MB \cite{xu2021}). No contradictions; variations stem from task-specific architectures (e.g., YOLO vs. ResNet). Hardware not ESP32-specific. Confidence: Strong (consistent performance gains with quantitative metrics).

\subsubsection{Microcontroller Deployment and Performance on Constrained Hardware}
Deployments on low-power MCUs like MAX78000 achieve 180 fps, 196 µJ/inference, $>10^6$ MAC/cycle \cite{moosmann2023}; Arduino Portenta supports on-device CNN training \cite{aatab2023}; STM32/Apollo4b tested for quantized models \cite{moosmann2023}. No ESP32 deployments reported. Edge devices generally hit real-time FPS (16-36.7 \cite{zhao2020, xu2021}) with 89-93\% precision \cite{heredia2019}. Power/memory metrics sparse beyond TinyissimoYOLO. Confidence: Limited (few MCU-specific metrics; no ESP32).

\subsubsection{Surveillance Tasks and Edge-Cloud Hybrids}
Object detection (handgun/knife \cite{liu2023}, pedestrians \cite{xu2021}, humans \cite{nikouei2018sec}) and tracking dominate, with edge-cloud cooperation boosting AP to 80.7\% via verification \cite{xu2021, lyu2022}. Minor/face privacy tasks integrate classification pipelines \cite{yuan2020}. Wake Vision enables TinyML person detection benchmarks (6.6\% accuracy gain \cite{banbury2024}). No method conflicts; hybrids address edge accuracy limits. Confidence: Moderate (task-consistent, real-world applications).

\subsubsection{On-Device Training and Adaptation in TinyML}
Libraries enable CNN/autoencoder training on MCUs ($\sim$hundreds KB RAM \cite{roveri2022, aatab2023}), adapting to data drifts without off-device retraining. No performance metrics or surveillance ties reported. Confidence: Limited (functional validation only; general-purpose).

\subsection{Summary of Evidence}

\begin{landscape}
\footnotesize
\begin{longtable}{ p{3.5cm} p{4.5cm} p{3.5cm} p{2cm} p{3.5cm} p{4cm} }
\caption{Summary of Evidence} \label{tab:summary} \\
\toprule
\textbf{Theme} & \textbf{Key Finding} & \textbf{Population Applicability} & \textbf{Effect Direction} & \textbf{Confidence Level} & \textbf{Supporting Studies} \\
\midrule
\endfirsthead
\toprule
\textbf{Theme} & \textbf{Key Finding} & \textbf{Population Applicability} & \textbf{Effect Direction} & \textbf{Confidence Level} & \textbf{Supporting Studies} \\
\midrule
\endhead
\midrule
\multicolumn{6}{r}{\textit{Continued on next page}} \\
\endfoot
\bottomrule
\endlastfoot

Privacy-Preserving Local Processing & 92.1\% minor classification accuracy; chaotic scrambling feasible & Smart city/campus surveillance & Positive & Moderate (consistent findings with reasonable design quality) & Myneni et al. \cite{myneni2022}, Yuan et al. \cite{yuan2020}, Huang et al. \cite{huang2023} \\

CNN Quantization and Lightweight Optimizations & 180 fps, 196 µJ/inference; 0.91\% mAP gain; 4-bit $>$ state-of-art & TinyML edge inference & Positive & Strong (consistent across multiple studies with metrics) & Moosmann et al. \cite{moosmann2023}, Chin et al. \cite{chin2022}, Kim et al. \cite{kim2020} \\

Microcontroller Deployment and Performance & 16-180 fps on MCUs; no ESP32 & Low-power MCUs (no ESP32) & Positive & Limited (sparse MCU-specific metrics) & Zhao et al. \cite{zhao2020}, Moosmann et al. \cite{moosmann2023} \\

Surveillance Tasks and Edge-Cloud Hybrids & 76.7-80.7\% AP pedestrians; 92-93\% precision counting & Industrial/urban surveillance & Positive & Moderate (consistent real-world tasks) & Xu et al. \cite{xu2021}, Lyu et al. \cite{lyu2022} \\

On-Device Training and Adaptation & Functional CNN training on $\sim$hundreds KB RAM MCUs & General TinyML & Positive & Limited (no metrics/surveillance) & Aatab et al. \cite{aatab2023}, Roveri et al. \cite{roveri2022} \\

\end{longtable}
\end{landscape}

\normalsize

\section{Discussion}
\subsection{Principal Findings and Their Interpretation}
Quantization emerges as the cornerstone enabling TinyML surveillance, with 8-bit \cite{moosmann2023} and adaptive 4-bit methods \cite{chin2022} preserving accuracy under memory extremes ($<$0.5 MB), because they dynamically mitigate biased activation losses inherent to modern CNNs like YOLO/ResNet---explaining superior metrics (180 fps, 196 µJ/inference) on MAX78000 versus unoptimized baselines. This pattern, visible only across studies, reveals why depthwise separable convolutions \cite{zhao2020, xu2021} synergize with quantization: they slash parameters pre-compression, amplifying edge feasibility for detection (89\% precision at 16 fps). Privacy gains stem mechanistically from local feature extraction (e.g., cascaded face pipelines at 92.1\% accuracy \cite{yuan2020}) and chaotic scrambling \cite{fitwi2021}, blocking identifiable data leaks without transmission---a direct counter to cloud vulnerabilities. Confidence is high in quantization efficacy (strong evidence, metrics consistent), moderate for privacy (architectural consistency sans universal metrics), and low for ESP32 (zero direct evidence). On-device training \cite{aatab2023, roveri2022} hints at adaptive surveillance, but absent metrics limit causal claims to functional plausibility. Collectively, synthesis underscores edge-local processing as viable, advancing beyond siloed reports by linking optimizations to surveillance outcomes.

\subsection{Comparison with Existing Literature and Resolution of Contradictions}
Findings align with broader edge literature on latency/bandwidth wins (e.g., edge at 92\% precision rivals server's 93\% \cite{heredia2019}), mechanistically because local CNNs eliminate upload delays, robustly confirmed across urban/industrial contexts. No major contradictions emerge, but sparse metrics (e.g., "not reported" in 12/20 studies) versus standouts like 180 fps \cite{moosmann2023} reflect methodological heterogeneity: MCU-specific papers quantify rigorously, while architecture-focused ones prioritize feasibility. This gap, not population differences, explains variability---earlier works (2018-2020) emphasize concepts \cite{nikouei2018sec}, later ones (2022-2024) add metrics via matured TinyML tools. Publication bias risk is moderate, as positive edge benefits dominate, potentially underreporting deployment failures; however, limitations like accuracy trade-offs are acknowledged \cite{moosmann2023}. Compared to centralized baselines, edge consistently reduces costs \cite{zhao2020}, with hybrids resolving residual edge limits via cloud verification \cite{xu2021}---a progression from pure edge, enhancing reliability without contradicting core claims.

\subsection{Practical Implications}
Privacy-sensitive deployments (smart cities, campuses) benefit most from local quantization/anonymization, enabling scalable CCTV without data aggregation for facilities lacking bandwidth \cite{myneni2022, yuan2020}; operators should prioritize MAX78000/STM32-like MCUs for 100+ fps detection under milliwatt power. Industrial sites (coal mines, airports) gain real-time safety via hybrids like FL-YOLO/Edge-YOLO (36.7-80.7\% AP \cite{xu2021, lyu2022}), advising edge for initial triage, cloud for verification. Clinically, minor protection pipelines suit schools/childcare \cite{yuan2020}, with 92.1\% accuracy guiding anonymization protocols. Public health favors edge to minimize breach risks in populated areas \cite{huang2023}. No threshold issues, as benefits scale continuously with edge adoption. Caveats: implications proxy to ESP32 via similar MCUs; untested ESP32 risks unknown Wi-Fi/memory interactions.

\subsection{Strengths and Limitations}
Strengths include comprehensive hybrid search capturing evolving TinyML/surveillance intersection and thematic synthesis revealing optimization-surveillance links. Included studies offer real-world prototypes (e.g., coal mine/airport). Limitations of studies: predominant conceptual work without uniform metrics/hardware (no ESP32), partial surveillance focus, and pre-2019 inclusions despite criteria. Review limitations: abstract-based screening risks missing details, no formal bias assessment, extraction prioritized structured data potentially overlooking nuances.

\section{Gaps and Future Directions}
Synthesis reveals critical gaps: zero ESP32 deployments despite its surveillance suitability (Wi-Fi, low cost), leaving microcontroller performance unbenchmarked for the exact RQ hardware. Quantitative metrics absent in most privacy/architecture papers hinder comparability, with only isolated FPS/energy (e.g., 180 fps \cite{moosmann2023}). Surveillance underrepresented in TinyML training works \cite{aatab2023, roveri2022}, missing adaptive detection evidence. Contradictions in metric sparsity unresolved due to design variation. Future studies must validate quantized CNNs (8/4-bit) on ESP32 for object detection, reporting FPS/power/accuracy on custom datasets like Wake Vision \cite{banbury2024}. Methodological advances: standardized TinyML benchmarks across ESP32/STM32/MAX78000, on-device QAT for surveillance drifts, and hybrid privacy metrics. Target underrepresented contexts: battery-powered outdoor nodes, multi-camera scalability.

\section{Conclusion}
Edge AI surveillance via TinyML with CNN quantization supports privacy-preserving local processing on low-power microcontrollers, achieving up to 180 fps and 196 µJ/inference on MAX78000 \cite{moosmann2023}, 16 fps at 89\% precision via depthwise separations \cite{zhao2020}, and 92.1\% minor detection for anonymization \cite{yuan2020}, though no studies deploy on ESP32---evidence proxies from similar MCUs. Quantization reliably enables real-time tasks (76.7-80.7\% AP pedestrians \cite{xu2021}), with edge outperforming clouds in latency/privacy via local feature extraction and scrambling \cite{fitwi2021}. This holds moderate-to-strong confidence for general constrained hardware, limited by ESP32 absence and metric gaps. The pivotal uncertainty is ESP32-specific feasibility, as its $\sim$520 KB SRAM/Wi-Fi may alter quantization trade-offs unseen in reviewed platforms. Advancing ESP32 benchmarks will unlock widespread, ESP32-native surveillance, transforming edge AI from prototype to production.

\begin{thebibliography}{99}

\bibitem{aatab2023} Aatab, S., \& Freitag, F. (2023). Towards a Library of Deep Neural Networks for Experimenting with on-Device Training on Microcontrollers. \textit{2023 IEEE 9th World Forum on Internet of Things (WF-IoT)}. IEEE. %%%

\bibitem{banbury2024} Banbury, C. R., et al. (2024). Wake Vision: A Tailored Dataset and Benchmark Suite for TinyML Computer Vision Applications. %%%

\bibitem{chin2022} Chin, H.-H., Tsay, R.-S., \& Wu, H.-I. (2022). An Adaptive High-Performance Quantization Approach for Resource-Constrained CNN Inference. \textit{2022 IEEE 4th International Conference on Artificial Intelligence Circuits and Systems (AICAS)}. IEEE. %%%

\bibitem{fagbohungbe2022} Fagbohungbe, O., et al. (2022). Efficient Privacy Preserving Edge Intelligent Computing Framework for Image Classification in IoT. \textit{IEEE Transactions on Emerging Topics in Computational Intelligence}. %%%

\bibitem{fitwi2021} Fitwi, A., Chen, Y., \& Zhu, S. (2021). Lightweight Frame Scrambling Mechanisms for End-to-End Privacy in Edge Smart Surveillance. IEEE. %%%

\bibitem{gheisari2019} Gheisari, M., et al. (2019). ECA: An Edge Computing Architecture for Privacy-Preserving in IoT-Based Smart City. \textit{IEEE Access}. %%%

\bibitem{hailesellasie2019} Hailesellasie, M. T., \& Hasan, S. R. (2019). MulNet: A Flexible CNN Processor With Higher Resource Utilization Efficiency for Constrained Devices. \textit{IEEE Access}. %%%

\bibitem{heredia2019} Heredia, A., \& Barros-Gavilanes, G. (2019). Edge vs. cloud computing: where to do image processing for surveillance? \textit{Artificial Intelligence and Machine Learning in Defense Applications}. SPIE. %%%

\bibitem{huang2023} Huang, A. Y. C., et al. (2023). Semantic Privacy-Preserving for Video Surveillance Services on the Edge. \textit{Proceedings of the Eighth ACM/IEEE Symposium on Edge Computing}. %%%

\bibitem{kim2020} Kim, S., \& Kim, H. (2020). Mixture of Deterministic and Stochastic Quantization Schemes for Lightweight CNN. \textit{2020 International SoC Design Conference (ISOCC)}. IEEE. %%%

\bibitem{liu2023} Liu, H., \& HU, Z. (2023). An Edge Computing-based Handgun and Knife Detection Method in IoT Video Surveillance Systems. \textit{International Journal of Advanced Computer Science and Applications}. %%%

\bibitem{lyu2022} Lyu, Z., \& Luo, J. (2022). A Surveillance Video Real-Time Object Detection System Based on Edge-Cloud Cooperation in Airport Apron. \textit{Applied Sciences}. %%%

\bibitem{moosmann2023} Moosmann, J., et al. (2023). TinyissimoYOLO: A Quantized, Low-Memory Footprint, TinyML Object Detection Network for Low Power Microcontrollers. \textit{2023 IEEE 5th International Conference on Artificial Intelligence Circuits and Systems (AICAS)}. %%%

\bibitem{myneni2022} Myneni, S., et al. (2022). SCVS: On AI and Edge Clouds Enabled Privacy-preserved Smart-city Video Surveillance Services. \textit{ACM Transactions on Internet of Things}. %%%

\bibitem{nikouei2018sec} Nikouei, S. Y., Chen, Y., \& Faughnan, T. R. (2018). Smart Surveillance as an Edge Service for Real-Time Human Detection and Tracking. \textit{2018 IEEE/ACM Symposium on Edge Computing (SEC)}. %%%

\bibitem{roveri2022} Roveri, P. M., \& Pavan, M. (2022). TinyML on-device neural network training. %%%

\bibitem{xu2021} Xu, Z., Li, J., \& Zhang, M. (2021). A Surveillance Video Real-Time Analysis System Based on Edge-Cloud and FL-YOLO Cooperation in Coal Mine. \textit{IEEE Access}. %%%

\bibitem{yuan2020} Yuan, M., et al. (2020). Minor Privacy Protection Through Real-time Video Processing at the Edge. \textit{2020 29th International Conference on Computer Communications and Networks (ICCCN)}. %%%

\bibitem{zhao2020} Zhao, Y., Yin, Y., \& Gui, G. (2020). Lightweight Deep Learning Based Intelligent Edge Surveillance Techniques. \textit{IEEE Transactions on Cognitive Communications and Networking}. %%%

\end{thebibliography}

\end{document}
